    \documentclass[11pt,a4paper]{article}
    \usepackage[utf8]{inputenc}
    \usepackage[T1]{fontenc}
    \usepackage{amsmath,amssymb}
    \usepackage{booktabs}
    \usepackage{geometry}
    \usepackage{hyperref}
    \usepackage{siunitx}
    \usepackage{pifont}
    \geometry{margin=2.5cm}

    \title{Definizioni e utili}
    \author{Marco Chiodo}
    \date{Definizioni e utili}

    \begin{document}
    \maketitle


    \begin{itemize}
        \item \textbf{Tempo di Risposta}: Tempo necessario ad un sistema per reagire ad una richiesta.
        \item \textbf{Throughput}: Numero di richieste evase per unità di tempo.
        \item \textbf{Availability}: Frazione di tempo durante la quale il sistema è attivo e accessibile per i suoi utenti.
        \item \textbf{Scalability}: Un sistema è detto \textit{scalabile} se le sue prestazioni non degradano significativamente al crescere del numero dei suoi utenti.
        \item \textbf{Thrashing} $\Rightarrow$ \textbf{No Admission Control} $\Rightarrow$ Non respinge le richieste in eccesso.
        \item \textbf{No Thrashing} $\Rightarrow$ \textbf{Admission Control} $\Rightarrow$ Respinge le richieste in eccesso.
        \item \textbf{Modelli Simulativi}: Basati su un programma che simula i diversi comportamenti del sistema.
        \item \textbf{Modelli Analitici}: Costituiti da un insieme di formule e/o algoritmi che forniscono i valori delle misure di prestazioni desiderate in funzione del workload.
        \item \textbf{Service Demand}: Tempo totale medio speso da una transazione in una data risorsa.
    \end{itemize}

    \section{QN}
    \subsection{Classi di Workload Aperte}
    Una classe di workload è detta aperta quando:
    \begin{itemize}
        \item L'intensità di workload è specificato tramite un tasso di arrivi.
        \item Il numero di clienti nel sistema è potenzialmente illimitato.
        \item Il throughput è fornito come parametro di input (in equilibrio è pari al tasso di arrivi)
    \end{itemize}

    \subsection{Classi di Workload Chiuse}
    Una classe di workload è detta chiusa se:
    \begin{itemize}
        \item L'intensità del workload è specificata dalla popolazione interna del sistema.
        \item Il numero di clienti del sistema è limitato e conosciuto.
        \item Il throughput è un parametro di output. Si ottiene risolvendo il modello QN.
    \end{itemize}

    \subsection{TIpi di Risorse}
    Possiamo usare nei modelli QN 3 tipi di risorse:
    \begin{itemize}
    \item \textbf{Load Independent}: Il tasso di servizio non dipende dal carico o dal numero di richieste in coda (Si pensi a un singolo processore o a un disco fisso: il tempo fisico necessario per eseguire una determinata istruzione o leggere un blocco di dati rimane costante, indipendentemente da quante altre transazioni stiano aspettando il loro turno nella coda).
    \item \textbf{Load Dependent}: Il tasso di servizio è funzione del numero di richieste in coda (Si pensi ad una lan condivisa da più utenti, più utenti sono connessi simultaneamente più la larghezza di banda per utente si ridurrà).
    \item \textbf{Delay}: Non ci sta una coda (Si pensi al "think time" di un utente umano: il tempo che l'utente impiega per leggere i contenuti di una pagina web prima di cliccare su un link successivo introduce un ritardo puro, senza che le richieste debbano accodarsi per competere per una risorsa. È utile anche per modellare componenti dedicati in modo esclusivo a una richiesta).
    \end{itemize}

    \subsection{Notazione di una QN}
    \begin{itemize}
        \item \textbf{$K$}: Numero di code in una QN.
        \item \textbf{$R$}: Numero di classi in una QN.
        \item \textbf{$\overrightarrow{\lambda} = (\lambda_1,...,\lambda_r)$}: Vetore dei tassi di arrivo per una QN aperta multi-classe.
        \item \textbf{$\overrightarrow{N} = (N_1,...,N_r)$}: Vetore della popolazione per una QN chiusa multi-classe.
        \item \textbf{$D_{i,r}$}: Service Demand dei clienti di classe $r$ presso la coda $i$.
        \item \textbf{$\text{Prior(r)}$}: Priorità della classe $r$.
        \item \textbf{$X_{0,r}$}: Throughput di richieste di classe $r$.
        \item \textbf{$R_{0,r}$}: Tempo di risposta per richieste di classe $r$.
    \end{itemize}

    \section{Parametri}
    \subsection{Variabili Operazionali}
    Solitamente misurabili direttamente:
    \begin{itemize}
        \item \textbf{T}: Lunghzza del periodo di osservazione.
        \item \textbf{K}: Numero di risorse nel sistema.
        \item \textbf{$B_i$}: Tempo in cui la risorsa $i$ è occupata (busy) nel periodo di osservazione $T$.
        \item \textbf{$A_i$}: Numero totale di richieste di servizio (o arrivi) alla risorsa i nel periodo di osservazione $T$.
        \item \textbf{$A_0$}: Numero totale di richieste arrivate al sistema nel periodo di osservazione $T$.
        \item \textbf{$C_i$}: Numero totale di richieste evase dalla risorsa i nel periodo di osservazione $T$.
        \item \textbf{$C_0$}: Numero totale di richieste evase dal sistema nel periodo di osservazione $T$.
    \end{itemize}

    \subsection{Metriche Derivate}
    \begin{itemize}
        \item \textbf{$S_i = \frac{B_i}{C_i}$}: Tempo medio di servizio per la risorsa $i$.
        \item \textbf{$U_i = \frac{B_i}{T}$}: Utilizzazione della risorsa $i$.
        \item \textbf{$X_i = \frac{C_i}{T}$}: Throughput della risorsa $i$.
        \item \textbf{$\lambda_i = \frac{A_i}{T}$}: Tasso degli arrivi alla risorsa $i$.
        \item \textbf{$X_0 = \frac{C_0}{T}$}: Throughput del sistema.
    \end{itemize}

    \section{Leggi Operazionali}
    
    \subsection{Utilization Law}
    Definiamo l'utilizzazione come il tempo medio di servizio per il throughput: $U_i = S_i\cdot X_i$ nel caso dove la risorsa $i$ è in equilibrio allora abbiamo che $\lambda_i = X_i$, quindi $U_I = S_i\cdot \lambda_i$.
    \begin{itemize}
        \item In caso la risosa $i$ disponga di $m$ server (cpu multi-core) allora: $U_i = \frac{S_i\cdot X_i}{m}$
    \end{itemize}

    \subsection{Service Demand Law}
    Il Service Demand $D_i$ indica il tempo totale medio speso da una richiesta nella risorsa $i$, esso è pari alla somma dei singoli tempi di visita, fino ad evasione della richiesta (una richiesta può visitare una risorsa più volte prima di essere evasa)
    Posso calcolarlo come:
    \begin{equation}
        D_i = \frac{B_i}{C_0} = \frac{U_i\cdot T}{C_0} = \frac{U_i}{\frac{C_0}{T}} = \frac{U_i}{X_0}
    \end{equation}
    Oppure più intuitivamente come il tempo di servizio medio della risorsa $i$ moltiplicato al numero di visite medio della risorsa $i$: $D_i = S_i\cdot V_i$

    \subsection{Forced Flow Law}
    Ci serve lper legare il throughput di una risorsa $i$, $X_i$ a quello del sistema $X_0$:
    \begin{equation}
        X_i = X_0 \cdot V_i
    \end{equation}
    
    \subsection{Little's Law}
    Ci dice il numero medio di clienti in un centro $i$, $N_i$, a patto che non vengano creati o distrutti clienti.
    \begin{equation}
        N_i = X_i\cdot R_i
    \end{equation}
    Dove $R_i$ è il tempo di risposta medio del centro $i$ inclusivo del tempo di coda\\
    Possiamo anche calcolare il nuemro delle sole richeste in coda: $N_i^W = X_i\cdot W_i$, dove $W_i$ è il tempo di coda medio al centro $i$.
    Possiamo anche calcolare il numero medio di clienti in un centro come il tempo medio trascorso nel centro moltiplicato al tasso di espulsione.

    \subsection{Interactive Response Time Law}
    Simulativi $Z$ il thinktime e diciamo che ci sono $M$ utenti, allora siano $\bar{M}$ il numero di utenti in fase di think e $\bar{N}$ il numero di utenti in attesa di risposta, quindi $M = \bar{M} + \bar{N}$.\\
    Per la legge di Little, abbiamo anche che:
    \begin{equation}
    \bar{M} = X_0\cdot Z\qquad \bar{N} = X_0\cdot R\qquad M = X_0(Z+R)\qquad R = \frac{M}{X_0} - Z
    \end{equation}


    \end{document}